\documentclass[12pt]{article}

\usepackage{amsmath, amssymb}
\usepackage{geometry}
\geometry{margin=1in}

\title{Lecture Scribe}
\author{Course: CSE400 -- Fundamentals of Probability in Computing}
\date{Topic: Joint Random Variables and Joint Distributions \\ Lectures: L15--L18}

\begin{document}

\maketitle

\section{Definitions and Notation}

\subsection{Joint Random Variables}

Let $X$ and $Y$ be two random variables defined on the same probability space.  
The pair $(X,Y)$ is called a \textbf{joint random variable}.

\subsection{Joint Probability Mass Function (Discrete Case)}

For discrete random variables $X$ and $Y$, the joint PMF is defined as

\[
p_{X,Y}(x,y) = P(X = x, Y = y)
\]

This gives the probability that the random variables $X$ and $Y$ simultaneously take values $x$ and $y$.

\subsection{Joint Probability Density Function (Continuous Case)}

For continuous random variables $X$ and $Y$, the joint PDF is defined as

\[
f_{X,Y}(x,y)
\]

such that

\[
P((X,Y) \in A) = \int\int_A f_{X,Y}(x,y)\, dx\, dy
\]

for any region $A$ in the plane.

\subsection{Joint Cumulative Distribution Function}

The joint CDF of two random variables $X$ and $Y$ is defined as

\[
F_{X,Y}(x,y) = P(X \le x, Y \le y)
\]

\section{Assumptions}

The following assumptions are used in the lecture:

\begin{itemize}
\item Random variables $X$ and $Y$ are defined on the same probability space.
\item For the discrete case, probabilities are represented using a joint PMF table.
\item For the continuous case, probabilities are computed using double integration over a region.
\item The joint PDF satisfies non-negativity and normalization conditions.
\item The joint CDF represents cumulative probability up to the point $(x,y)$.
\end{itemize}

\section{Theorems / Results}

\subsection{Properties of the Joint PMF}

\textbf{Non-negativity}

\[
p_{X,Y}(x,y) \ge 0
\]

\textbf{Normalization}

\[
\sum_x \sum_y p_{X,Y}(x,y) = 1
\]

\subsection{Properties of the Joint PDF}

\textbf{Non-negativity}

\[
f_{X,Y}(x,y) \ge 0
\]

\textbf{Normalization}

\[
\int_{-\infty}^{\infty}\int_{-\infty}^{\infty} f_{X,Y}(x,y)\, dx\, dy = 1
\]

\subsection{Joint CDF Expression (Continuous Case)}

\[
F_{X,Y}(x,y) =
\int_{-\infty}^{x}\int_{-\infty}^{y}
f_{X,Y}(u,v)\, du\, dv
\]

\subsection{Joint CDF Expression (Discrete Case)}

\[
F_{X,Y}(x,y) =
\sum_{u \le x}\sum_{v \le y}
P(X = u, Y = v)
\]

\subsection{Properties of the Joint CDF}

\[
0 \le F_{X,Y}(x,y) \le 1
\]

$F_{X,Y}(x,y)$ is non-decreasing in both arguments.

Limits:

\[
\lim_{x \to -\infty} F_{X,Y}(x,y) = 0
\]

\[
\lim_{y \to -\infty} F_{X,Y}(x,y) = 0
\]

\[
\lim_{x,y \to \infty} F_{X,Y}(x,y) = 1
\]

\section{Proofs / Derivations}

\subsection{Derivation of Joint CDF from Joint PDF}

Starting from the definition:

\[
F_{X,Y}(x,y) = P(X \le x, Y \le y)
\]

For continuous random variables:

\[
F_{X,Y}(x,y) =
\int_{-\infty}^{x}\int_{-\infty}^{y}
f_{X,Y}(u,v)\, du\, dv
\]

This follows from integrating the joint PDF over the region

\[
(-\infty, x] \times (-\infty, y]
\]

\subsection{Probability over a Region (Continuous Case)}

If $A$ is a region in the plane:

\[
P((X,Y) \in A) =
\int\int_A f_{X,Y}(x,y)\, dx\, dy
\]

This represents the probability that the pair $(X,Y)$ lies in region $A$.

\section{Worked Examples}

\subsection{Example 1: Rolling Two Dice}

Let

\[
X = \text{value on the first die}
\]

\[
Y = \text{value on the second die}
\]

Each die can take values:

\[
1,2,3,4,5,6
\]

\subsection{Joint PMF}

Since the dice are fair:

\[
p_{X,Y}(x,y) = \frac{1}{36}
\]

for

\[
x,y \in \{1,2,3,4,5,6\}
\]

\subsection{Event Example}

Consider the event

\[
X + Y = 7
\]

Possible outcomes:

\[
(1,6),(2,5),(3,4),(4,3),(5,2),(6,1)
\]

Total outcomes = $6$

Therefore

\[
P(X+Y=7) = \frac{6}{36} = \frac{1}{6}
\]

\subsection{Continuous Case Example}

Let $f_{X,Y}(x,y)$ be defined over a region in the plane.

The probability over a region $A$ is obtained using

\[
P((X,Y) \in A) =
\int\int_A f_{X,Y}(x,y)\, dx\, dy
\]

The integration limits correspond to the boundaries of the region shown in the lecture.

\section{Conclusion}

The lecture introduced the concept of joint random variables and their probability distributions.

The following were covered:

\begin{itemize}
\item Joint PMF for discrete random variables
\item Joint PDF for continuous random variables
\item Geometric interpretation of joint distributions
\item Definition of the joint CDF
\item Formulas for computing joint CDF in discrete and continuous cases
\item Properties of the joint CDF
\item Example of rolling two dice using a joint PMF table
\end{itemize}

\end{document}